\begin{reviewed}

	\section{Literature Review Methodology}

	The first phase of the methodology involves constructing a search string composed of relevant keywords to retrieve articles from target databases. Due to the rapid growth of the \gls{ai} safety domain and the occasional overlap of topics, several iterations are performed to find the most suitable search string. This iterative process ensures the retrieval of literature that adequately represents the area, while a negative search string filters out works from specific domains outside the scope of this research. The detailed review protocol, including the final search string and the specific databases used, is presented in Chapter \ref{chap:literature_review}.

	Furthermore, based on the work of J. Webster and R. T. Watson on writing a literature review \cite{webster_analyzing_2002}, this methodology utilizes a concept matrix to assist in discerning the different articles after the abstract screening phase. The final matrix is available in Appendix \ref{sec:search_string__articles_slr}.

\begin{toreview}

	The seven categories used in this Concept Matrix were:
	\begin{itemize}
		\item Value Alignment
		\item Auditing Safety \& oversight
		\item Explainability / Interpretability
		\item Bias \& Fairness
		\item Human \& Societal Impacts
		\item Governance, regulation \& Policy
		\item Liability \& Ethics
	\end{itemize}
	
	These do not emerge from a single taxonomy. Instead, they combine two different sources that each cover part of the field. For the technical categories (Value Alignment, Auditing Safety \& Oversight, Explainability/Interpretability), the four-pillar structure from \textit{Unsolved Problems in ML Safety} \cite{hendrycks2021unsolved} was used as a starting point, since it organizes AI safety around Robustness, Monitoring, Alignment, and Systemic Safety. This did not, however, cover the societal and policy side of the field, so the Domain Taxonomy from the MIT AI Risk Repository \cite{Slattery_2026} was also consulted, since it aggregates 43+ existing AI risk taxonomies into seven domains. These were distilled into \textit{Bias \& Fairness} and \textit{Human \& Societal Impacts}. The categories \textit{Governance, regulation \& Policy} and \textit{Liability \& Ethics} were then added to account for the large amount of policy oriented literature.

\end{toreview}

	Additionally, the methodology employs snowballing to integrate other relevant articles and gray literature into the review \cite{wohlin_guidelines_2014}.

\end{reviewed}
